\documentclass[11pt,a4paper]{article}
\usepackage[T1]{fontenc}
\usepackage[utf8]{inputenc}
\usepackage{lmodern}
\usepackage[margin=27mm,headheight=14pt]{geometry}
\usepackage{amsmath,amssymb,amsthm,mathtools}
\usepackage{booktabs,tabularx,longtable,array}
\usepackage{enumitem}
\usepackage{microtype}
\usepackage{xcolor}
\usepackage{fancyhdr}
\usepackage{hyperref}
\usepackage{bookmark}
\usepackage{xurl}
\definecolor{ink}{HTML}{172D43}
\definecolor{accent}{HTML}{1A6582}
\hypersetup{colorlinks=true,linkcolor=accent,citecolor=accent,urlcolor=accent,
 pdftitle={Two order tests instead of primitive polynomial labels},
 pdfauthor={OpenAI ChatGPT},
 pdfsubject={A source-audited local simplification of the ProveIt theta-lift proof of Dirichlet GRH}}
\pagestyle{fancy}
\fancyhf{}
\fancyhead[L]{\small Two order tests instead of primitive labels}
\fancyhead[R]{\small Technical proposal}
\fancyfoot[C]{\thepage}
\setlength{\parskip}{4pt}
\setlength{\parindent}{0pt}
\setlist{itemsep=3pt,topsep=5pt}
\newtheorem{theorem}{Theorem}[section]
\newtheorem{lemma}[theorem]{Lemma}
\newtheorem{proposition}[theorem]{Proposition}
\newtheorem{corollary}[theorem]{Corollary}
\theoremstyle{definition}
\newtheorem{definition}[theorem]{Definition}
\theoremstyle{remark}
\newtheorem{remark}[theorem]{Remark}
\newcommand{\F}{\mathbb F}
\newcommand{\N}{\mathbb N}
\newcommand{\ord}{\operatorname{ord}}
\newcommand{\Norm}{\operatorname{N}}
\newcommand{\good}{\mathcal S}
\newcommand{\labels}{\mathcal C}
\newcommand{\code}[1]{{\small\nolinkurl{#1}}}
\newcommand{\src}[2]{\href{https://github.com/VladimirReshetnikov/ProveIt/blob/89d12dd40424f1217fd818d84877b3da52396a3b/Analysis/GeneralizedRiemannHypothesis/Lean/grh/src/#1.lean}{#2}}
\newcommand{\guide}{\href{https://github.com/VladimirReshetnikov/ProveIt/blob/89d12dd40424f1217fd818d84877b3da52396a3b/Analysis/GeneralizedRiemannHypothesis/docs/guide/theta_lifts_and_dirichlet_zeros.tex}{human-readable guide}}
\newcommand{\proposal}[2]{\href{https://github.com/VladimirReshetnikov/ProveIt/tree/89d12dd40424f1217fd818d84877b3da52396a3b/Analysis/GeneralizedRiemannHypothesis/docs/proposals/#1}{#2}}
\begin{document}
\hypersetup{pageanchor=false}
\begin{titlepage}
\color{ink}
\vspace*{12mm}
{\large SOURCE REVIEW AND A CONCRETE SIMPLIFICATION\par}
\vspace{8mm}
{\huge\bfseries Two order tests instead of\\[3pt] primitive polynomial labels\par}
\vspace{7mm}
{\Large A local replacement in the theta-lift proof\\ of the Riemann hypothesis and Dirichlet GRH\par}
\vspace{10mm}
{\large Review of the supplied ProveIt development\par}
{\large 27 November 2026\par}
\vspace{12mm}
\begin{abstract}
\noindent
The finite-field primitive-polynomial condition in the supplied proof is stronger than its
consumers require. For rank \(r=3^a\), replace a root generating
\(\F_{q^r}^{\times}\) by a root \(z\) satisfying two order tests,
\(z^{(q^r-1)/2}\ne1\) and \(z^{(q^r-1)/3}\ne1\), and one explicit subgroup
exclusion, \(z^{3(q^{r/3}-1)}\ne1\).
These conditions preserve full Frobenius orbits after every prime-power operation,
survive the original scalar normalization, and give a common scalar power of sufficient
order after weakening one internal interface. The resulting irreducible labels number
exactly \((q^r-3q^{r/3}+2)/(3r)\).
Consequently one block of residue fields of equal cardinality has bad share at most
\(f\varepsilon+\exp(-f/(6r))\), replacing the source's
\(f\varepsilon+\exp(-f/[2r(2r\log q+1)])\).
This removes the analytic totient estimate from this branch.
The proposal is distinct from all five existing proposal directories.
Full mathematical proofs, source anchors, executable exact checks, and a checked
Lean support file accompany this report. Integration into the complete GRH development
remains to be done.
\end{abstract}
\vfill
\small
\textbf{Source snapshot.}
\code{VladimirReshetnikov/ProveIt}, commit\\
\code{89d12dd40424f1217fd818d84877b3da52396a3b}.\\[4pt]
\textbf{Scope of verification.}
Source-level dependency review and a proved local replacement.
This report does not claim an independent rebuild or certification of the complete
original RH/GRH project.
\end{titlepage}
\hypersetup{pageanchor=true}
\pagenumbering{roman}
\setcounter{tocdepth}{1}
\tableofcontents
\newpage
\pagenumbering{arabic}

\section{Recommendation and scope}

I recommend option (a): simplify the arithmetic input that defines the good finite-field
types. This changes a small, identifiable mathematical interface while leaving the
theta-lift, congruence, Selmer, and analytic framework available for reuse.
The main gain is an exact, constant-density count of usable roots, with a direct proof of
every property that the subsequent argument needs.

Throughout, \emph{GRH} refers to the Dirichlet-\(L\) statements exported by the supplied
project. Its \src{Statements}{Statements.lean} explicitly quantifies over all moduli and
Dirichlet characters; it does not state GRH for arbitrary automorphic \(L\)-functions.
Its \src{Main}{Main.lean} exports Mathlib's \code{RiemannHypothesis} through
\code{New.zeta_zeros_on_line New.grh}, as well as related zero-location and nonvanishing
statements. The source snapshot is pinned so that the proposal can be reviewed against
unchanging definitions.

The audit inspected the guide, every textual source in the five proposal directories,
and the source interfaces surrounding the candidate replacement. The checked-out
\code{src} tree contains 803 Lean files. This is not a line-by-line mathematical
rederivation of all those files, and I did not rerun their complete build.
The conclusions below distinguish three kinds of evidence:

\begin{enumerate}
\item \textbf{Mathematical replacement.} Sections~\ref{sec:theorem}--\ref{sec:counting}
give proofs of the new finite-field, nonresonance, and weighted-count assertions.
\item \textbf{Source compatibility.} Sections~\ref{sec:uses} and \ref{sec:integration}
identify the exact consumers of primitivity, including the one order-divisibility
field that must change.
\item \textbf{Machine checks.} The archive contains independent exact finite-group
checks and standalone Lean support. Their scope is narrower than a complete
integration; see Section~\ref{sec:verification}.
\end{enumerate}

In particular, the existing unimported \code{Challenge.lean} contains intentional
comparison placeholders. Their presence is not evidence of a hole in the import
closure of the exported theorem. Conversely, reading the exported declaration alone
does not independently certify its entire dependency closure.

\subsection{What improves, quantitatively}

Let \(q\) be the residue-field cardinality in one block, \(r\) the rank, \(f\) the number
of factors in that block, \(M\) the original tame-character order, and
\[
  \varepsilon=\frac{1+M+2r}{q},\qquad 1<r<q.
\]
The source's regular-representation projector calculation has two error terms:
the cost \(f\varepsilon\) of forbidden labels and the cost of having no sufficiently
strong witness among the \(f\) factors. Only the latter changes.

\begin{center}
\begin{tabularx}{\textwidth}{@{}lXX@{}}
\toprule
 & Inspected source & Proposed replacement\\
\midrule
Witness root & Generator of \(\F_{q^r}^{\times}\) & Two order tests and one subgroup exclusion\\
Witness polynomials & \(\varphi(q^r-1)/r\) & \((q^r-3q^{r/3}+2)/(3r)\)\\
Uniform usable exponent & \(1/[2r(2r\log q+1)]\) & At least \(1/(6r)\)\\
Sufficient \(f\) for tail \(e^{-L}\) & \(2r(2r\log q+1)L\) & \(6rL\)\\
\bottomrule
\end{tabularx}
\end{center}

The table compares sufficient bounds in the same counting argument.
It does not say that actual primitive-label density is as small as the source's lower
bound, or that the complete proof becomes faster by the ratio of the two bounds.
For fixed \(r,q\), the exact new exponent in \eqref{eq:uexact} is better than the
simple \(1/(6r)\) estimate.

\section{Where this sits in the original argument}
\label{sec:architecture}

The source organizes the proof around a comparison between an upper bound obtained
from short integral forms and a lower bound obtained from a corrected chain
covolume. The following description tracks the mathematical interfaces relevant
to the proposed change rather than enumerating every supporting theory.

\subsection{The common inequality}

The theorem \code{star_inequality} in
\src{New/Reduction/MainInequality}{New/Reduction/MainInequality.lean}
combines three inputs:

\begin{enumerate}
\item short integral source forms, with their index and height costs controlled;
\item sufficiently many translates to occupy almost all of the required target;
\item a lower bound for the target's corrected covolume, obtained from chains of
theta lifts and arithmetic control of their projection denominators.
\end{enumerate}

In the source's notation, the resulting comparison has the form
\[
      m\,\mathcal R \le f a(\log b)^5 .
\]
Here \(\mathcal R\) is the relevant difference of smoothed prime remainders, \(m\)
is the chain length, and the remaining letters encode the admissible scale.
The exact statement, including all data hypotheses, remains the source theorem.
The present proposal changes one input to the good target projection; it does not
replace this comparison.

The source then chooses admissible data to derive the eventual one-sided bound
\[
             R_{q_0,d}(x)\le x^{-1/2+\epsilon}
\]
for a fixed auxiliary quadratic twist \(d\).
The zero criterion chooses a twist that leaves a hypothetical zero visible and
uses a Landau-type sign argument. The functional equation gives the other half
of the critical strip. See
\src{New/Reduction/OneSidedBound}{OneSidedBound.lean},
\src{New/Reduction/Zeros/Criterion}{Zeros/Criterion.lean}, and
\src{New/Main}{New/Main.lean}.

\subsection{Why local labels matter to both sides}

At a split level place with residue field \(\F_q\), a label is a monic squarefree
polynomial of degree \(r\) with nonzero constant term. The finite-group
representation theory associates an idempotent to such a label. A tuple of labels
defines a tensor idempotent; summing over good tuples defines the chosen target
projection.

There are two distinct requirements on a good tuple:

\begin{enumerate}
\item Every entry must be \emph{allowed}: it has no root in the base field whose
\(M\)-th power is \(1\), in addition to being a label. This is what annihilates the
weighted hyperplane averages used to eliminate unwanted contributions.
\item At least one entry in each of the two level-prime blocks must provide a
strong witness. The current definition asks for a primitive polynomial.
\end{enumerate}

Irreducibility of the witness supplies the cuspidal local test needed for trace
comparison across signatures. The same witness subsequently provides an integral
Galois-image statement: after removing the primary part at any coefficient
prime, its Frobenius orbit still has size \(r\). Distinct residual eigenvalues give
unit Vandermonde determinants, so the tame pair spans the full matrix algebra
integrally. Finally, a common power of those eigenvalues must have enough order
to stay separated from \(1\) after certain bounded powers.
These roles are explained in the \guide{} and implemented in
\src{New/UpperBound/Objects/LabelsAndIndexTypes}{LabelsAndIndexTypes.lean},
\src{New/Spanning/Multiplicities/EqualMultiplicities}{EqualMultiplicities.lean},
and \src{New/Congruences/GaloisRepresentations/RowsAndColumnsFromLabels}{RowsAndColumnsFromLabels.lean}.

This is why simply replacing \emph{primitive} by \emph{irreducible} would be invalid.
It would preserve the trace-comparison input while losing the arithmetic one.
The replacement below explicitly preserves both.

\section{Non-overlap with the existing proposals}

The exclusion audit covers all 20 textual proposal sources, including subsidiary
TeX files, READMEs, implementation plans, a Lean supplement, and future-work sections.
The archive records their paths and SHA-256 hashes in
\code{provenance/proposal_inventory.csv}.

\begin{center}
\begin{tabularx}{\textwidth}{@{}p{.31\textwidth}X@{}}
\toprule
Existing directory & Proposals treated as already available\\
\midrule
\proposal{centered_theta_lift_simplification}{Centered theta-lift simplification}
& Centered covariance, adaptive spanning, actual completion cost, random
projection products; future rank--height budgets and partial-volume approaches.\\
\proposal{proof_simplifications}{Proof simplifications}
& Direct spanning on the ramified source, Fourier or large-sieve twist selection,
and direct far-cutoff Perron estimates.\\
\proposal{rh-simplifications}{RH simplifications}
& Sharper rank recurrence and its abstract obstruction, quadratic large sieve,
finite torsion interpolation, and endpoint simplifications.\\
\proposal{rh_sparse_theta_proof}{Sparse theta proof}
& Reciprocal convexity, sparse cutoffs, local averaging, Gaussian twist
selection, favorable auxiliary-prime families, and shared blockwise data.\\
\proposal{shorter-theta-lift-proof}{Shorter theta-lift proof}
& A combined shorter route through twist selection, spanning, prime
asymptotics, and cutoff interpolation.\\
\bottomrule
\end{tabularx}
\end{center}

None of these changes the primitive-polynomial witness condition, replaces its
finite-field count, or weakens the common-root-order interface below. In particular,
favorable auxiliary \emph{quadratic} primes \(d\) are different from the residue-field
labels at the two already chosen level primes.

The proposal therefore remains useful if all the existing proposals are implemented:
their spanning and analytic improvements still need sufficiently many admissible
local types and the Galois properties of those types. This is a claim of
non-overlap with the supplied proposal directory, not a claim that the finite-field
observations are unprecedented in the mathematical literature.

\section{The exact interface currently supplied by primitivity}
\label{sec:uses}

The source predicate \code{IsPrimitivePolynomial r c} packages a label, its
irreducibility, and a root of multiplicative order \(q^r-1\). A direct search finds
this predicate in exactly six source files, listed in
Section~\ref{sec:integration}. A related generic orbit lemma lives in
\code{Lib/LocalGL/Green/LabelCharacters.lean}.

The substantive consumer is the structure \code{IsCommonRootPower}. In mathematical
notation it asks for an exponent \(e\) and a scalar \(y\) such that:

\begin{enumerate}[label=(\roman*)]
\item every root \(\xi\) of the normalized type polynomial has \(\xi^e=y\);
\item every root is fixed by the \(q^r\)-power map;
\item for each prime \(\ell\) under consideration and every \(u\ge0\), the
conjugates \((\xi^{q^a})^{\ell^u}\), \(0\le a<r\), remain distinct;
\item if \(y^x=1\), then \(M\mid jx\), where \(j\) is one of the source's
bounded twist powers.
\end{enumerate}

The new finite-field condition supplies (i)--(iii) directly. For (iv), it supplies
divisibility by a fixed \(2,3\)-smooth order that is sufficient for the sole
downstream use. Section~\ref{sec:order} proves that weakening explicitly.
It would be incorrect to retain the old field and claim that the new roots
always supply the full \(M\)-divisibility.

Scalar normalization matters as well. The original proof rescales roots by a
base-field scalar \(c\) with \(c^M=1\), using the congruence \(M^2\mid q-1\).
The new witness family is stable under this exact normalization. The actual \(M\)
and the congruence are retained.

\section{A replacement finite-field theorem}
\label{sec:theorem}

Let \(q\) be a prime power, let \(r=3^a\) with \(a\ge1\), let \(M\) be a
positive integer, and assume
\begin{equation}
       6\mid M,\qquad M^2\mid q-1 .
       \label{eq:hyp}
\end{equation}
These hold for the source parameters: \(M=q_2(s)=(q_0d)^{s+6}\),
\(18480\mid q_0\), and the level primes are \(1\) modulo \(M^2\).
See \src{New/Reduction/Data/FieldsAndPlaces}{FieldsAndPlaces.lean},
\src{New/Reduction/Data/AdmissibleData}{AdmissibleData.lean}, and
\src{New/Reduction/Data/Magnitudes}{Magnitudes.lean}.

Put
\begin{equation}
 E=\F_{q^r},\quad
 A=q^{r/3},\quad
 N=q^r-1=A^3-1,\quad
 h=3(A-1),\quad
 K=\frac{A^2+A+1}{3}.
 \label{eq:defs}
\end{equation}
Thus \(N=hK\). The integer \(K\) is well defined because \(q\equiv1\pmod3\).
The letter \(M\) here denotes the source's tame order \(q_2\), not a product of
level primes; this avoids a possible conflict with notation elsewhere in the guide.

\begin{definition}
Define the selected-root set
\begin{equation}
 \good_{q,r}=
 \left\{z\in E^\times:
      z^{N/2}\ne1,\quad z^{N/3}\ne1,\quad z^h\ne1
 \right\}.
 \label{eq:selection}
\end{equation}
Let \(\labels_{q,r}\) be the set of monic irreducible degree-\(r\) polynomials
over \(\F_q\) whose roots belong to \(\good_{q,r}\).
\end{definition}

The first two tests retain the full \(2\)-primary and \(3\)-primary components of
the root's order. The last excludes the subgroup
\[
 H=\{z\in E^\times:z^3\in\F_A^\times\}
   =\{z:z^{3(A-1)}=1\},
\]
which has order \(h\). The predicate itself does not depend on all prime factors
of \(M\).

\begin{theorem}[Replacement labels]
\label{thm:main}
Under \eqref{eq:hyp}, the following assertions hold.
\begin{enumerate}[label=(\alph*)]
\item For every \(z\in\good_{q,r}\), every prime \(\ell\), and every \(u\ge0\),
the element \(z^{\ell^u}\) has degree \(r\) over \(\F_q\).
\item The set \(\good_{q,r}\) is stable under Frobenius and under multiplication
by every \(c\in\F_q^\times\) with \(c^M=1\).
\item With
\[
             e=\frac{q^r-1}{q-1},\qquad y=z^e\in\F_q^\times,
\]
every Frobenius conjugate \(\xi\) of \(z\) satisfies \(\xi^e=y\).
For \(p=2,3\), the \(p\)-part of \(\ord(y)\) equals the \(p\)-part of \(q-1\).
In particular, if \(C\mid M\) and every prime factor of \(C\) is \(2\) or \(3\),
then \(C^2\mid\ord(y)\).
\item The exact counts are
\begin{align}
    |\good_{q,r}|&=\frac{q^r-3q^{r/3}+2}{3},\label{eq:rootcount}\\
    |\labels_{q,r}|&=\frac{q^r-3q^{r/3}+2}{3r}.\label{eq:labelcount}
\end{align}
Every primitive degree-\(r\) polynomial is in \(\labels_{q,r}\).
\end{enumerate}
\end{theorem}

\subsection{Elementary arithmetic of the excluded subgroup}

Write \(A=1+3v\). Then
\begin{equation}
       K=1+3v+3v^2\equiv1\pmod3 .
       \label{eq:kmod3}
\end{equation}
Since \(q\) is odd, \(A\) is odd; \(A^2+A+1\), and hence \(K\), is odd.
Thus \(\gcd(K,6)=1\). Also \(6\mid h\), and \(h\mid N\).
One can note the stronger congruence
\[
 K=1+(A-1)\frac{A+2}{3}\equiv1\pmod{q-1},
\]
because \(q-1\mid A-1\) and \(3\mid A+2\). Only \(\gcd(K,6)=1\)
is needed for the fixed two-test count.

\subsection{Full degree after every prime-power operation}

For \(z\in\good_{q,r}\), set \(w=z^{A-1}\).
Then
\begin{equation}
       w^{3K}=1,\qquad w^K=z^{N/3}\ne1,\qquad w^3=z^h\ne1.
       \label{eq:w}
\end{equation}
Let \(d_w=\ord(w)\). It divides \(3K\). Since \(3\nmid K\), the
condition \(w^K\ne1\) forces \(3\mid d_w\), with exactly one factor \(3\).
The condition \(w^3\ne1\) then forces \(d_w\) to have a prime factor other than
\(3\). Therefore \(d_w\) cannot divide a power of any single prime.
In particular,
\begin{equation}
                     w^{\ell^u}\ne1
                     \quad\text{for all primes \(\ell\) and \(u\ge0\)}.
                     \label{eq:nokill}
\end{equation}

Every proper intermediate field of \(E/\F_q\), where \(E=\F_{q^{3^a}}\), is
contained in its unique maximal proper intermediate field
\(\F_A=\F_{q^{3^{a-1}}}\). If \(z^{\ell^u}\) belonged to that
subfield, its \((A-1)\)-th power would be \(1\), contradicting \eqref{eq:nokill}.
It therefore has degree \(r\). This proves part (a), including ordinary
irreducibility by taking \(u=0\).

Equivalently, choose a generator \(g\) of \(E^\times\) and write \(z=g^k\).
The tests say
\begin{equation}
                     \gcd(k,6)=1,\qquad K\nmid k.
                     \label{eq:kselection}
\end{equation}
Membership of \(z^{\ell^u}\) in \(\F_A^\times\) would say
\(3K\mid k\ell^u\). For \(\ell\ne3\), the factor \(3\) already makes this
impossible. For \(\ell=3\), coprimality of \(K\) and \(3\) contradicts
\(K\nmid k\). This is a second, purely cyclic-group proof of the same assertion.

The condition needed in Lean is slightly more explicit than a degree assertion:
the map
\[
    a\longmapsto (z^{q^a})^{\ell^u},\qquad 0\le a<r
\]
must be injective. The degree assertion is exactly the assertion that the
Frobenius orbit of \(z^{\ell^u}\) has size \(r\), so it gives this injectivity.
Frobenius stability then gives the same assertion for every root of the minimal
polynomial, not just the root used to define it.

\subsection{Common power and the order that survives}

The norm exponent is \(e=1+q+\cdots+q^{r-1}=N/(q-1)\).
The element \(y=z^e\) is in \(\F_q^\times\). Consequently,
\[
       (z^{q^i})^e=y^{q^i}=y
       \qquad(0\le i<r).
\]
This is a common power of all roots; it does not require a primitive root.
It is also their product, so it agrees with the scalar chosen by the original
common-root-power proof. There is no sign adjustment.

For \(p=2,3\),
\[
             y^{(q-1)/p}=z^{N/p}\ne1.
\]
In a cyclic group of order \(q-1\), this means that the \(p\)-adic valuation of
\(\ord(y)\) is exactly that of \(q-1\).
Because \(M^2\mid q-1\), every \(2,3\)-smooth \(C\mid M\) has
\(C^2\mid\ord(y)\). This proves (c).

It is essential to restrict this last conclusion to \(2,3\)-smooth \(C\).
For example, at \(q=1801\), \(M=30\), a selected root can have norm order \(72\),
which has no factor \(5\). The statement \(M\mid\ord(y)\) would be false there.

\subsection{Stability under normalization}

Let \(c\in\F_q^\times\) with \(c^M=1\).
For \(p=2,3\), the hypotheses \(p\mid M\) and \(M^2\mid q-1\) imply
\[
                  M\mid\frac{q-1}{p}\mid\frac{N}{p}.
\]
Thus \(c^{N/p}=1\), and
\((cz)^{N/p}=z^{N/p}\ne1\).
Also \(q-1\mid A-1\), so \(c^h=1\) and
\((cz)^h=z^h\ne1\).
All three tests survive.

Frobenius preserves the tests because its action on \(E^\times\) is an
automorphism. This proves (b). Under polynomial root scaling, monicity and degree
are restored by the usual normalized polynomial, and irreducibility is preserved
by the invertible base-field substitution. These are exactly the facts needed by
the source's \code{scaleRoots} operation.

\subsection{Exact counting}

With \(z=g^k\), the first two tests exclude exponents divisible by \(2\) or \(3\).
Exactly \(N/3\) of the residues modulo \(N\) remain. The excluded subgroup \(H\)
consists of exponents \(k=Kt\), \(0\le t<h\).
Since \(\gcd(K,6)=1\), precisely \(h/3\) of these pass the first two tests.
Subtracting gives
\[
       |\good_{q,r}|=\frac{N-h}{3}
       =\frac{q^r-3q^{r/3}+2}{3}.
\]
By part (a), every Frobenius orbit in this set has length \(r\). Such orbits
correspond bijectively to monic irreducible degree-\(r\) polynomials, proving
\eqref{eq:labelcount}.

A generator \(g^k\) has \(\gcd(k,N)=1\). Since \(6\mid N\) and \(K>1\)
divides \(N\), it passes \eqref{eq:kselection}. Thus all original primitive labels
are retained. The inclusion is sometimes strict, but need not be strict for
every \(q,r\).

\section{Why two surviving order components are sufficient}
\label{sec:order}

The source's local scalar-order hypothesis is used to rule out an equality
\[
                       y^{\ell^b t}=1
\]
where \(\ell\) is the coefficient prime, \(1\le j\le m\),
\(1\le t\le4r+2\), and the scale hypotheses ensure
\begin{equation}
                         0<jt<2^s .
                         \label{eq:small}
\end{equation}
The relevant helper is \code{not_q2_dvd_mul_pow_mul} (with a subscript \(2\)
in the actual Lean name) in
\src{New/Congruences/GaloisRepresentations/RowsAndColumnsFromLabels}{RowsAndColumnsFromLabels.lean}.
Its original proof chooses \(2\) or \(5\) to obtain a prime different from \(\ell\).
Choosing \(2\) or \(3\) gives the following elementary replacement.

\begin{lemma}[Two-prime nonresonance]
\label{lem:nonresonance}
Let \(\ell\) be prime, and let \(j,t,s,b\) be nonnegative integers with
\(0<jt<2^s\). Then
\[
                     6^s\nmid j\ell^b t.
\]
In particular \(6^{s+6}\nmid j\ell^b t\).
\end{lemma}

\begin{proof}
Choose \(p=3\) when \(\ell=2\), and \(p=2\) otherwise. Then \(p\ne\ell\),
\(p\in\{2,3\}\), and \(p^s\mid6^s\).
If \(6^s\mid j\ell^b t\), coprimality of \(p^s\) and \(\ell^b\) would imply
\(p^s\mid jt\). Since \(jt>0\), this gives
\[
                     2^s\le p^s\le jt,
\]
contradicting \eqref{eq:small}. The final assertion follows from
\(6^s\mid6^{s+6}\).
\end{proof}

Now put
\begin{equation}
               C=6^{s+6},\qquad M=(q_0d)^{s+6}.
               \label{eq:orderC}
\end{equation}
Since \(6\mid q_0d\), we have \(C\mid M\).
Theorem~\ref{thm:main}(c) gives \(C^2\mid\ord(y)\), and therefore
\[
                       y^x=1\ \Longrightarrow\ C\mid x
                       \ \Longrightarrow\ C\mid jx .
\]
Combining this with Lemma~\ref{lem:nonresonance} yields the same exclusion
\(y^{\ell^b t}\ne1\) that the Galois-image argument consumes.

\subsection{The required source edit}

Change just the final field of the internal common-root-power interface from
\[
        y^x=1\ \Longrightarrow\ q_2(s)\mid jx
        \quad\hbox{to}\quad
        y^x=1\ \Longrightarrow\ 6^{s+6}\mid jx .
\]
Its only downstream field use in the pinned file is in the proof of the
scalar-inertia separation, where it is passed to the old nondivisibility helper.
Replace that call by the two-prime helper.
The ultimate \code{IsMainSummandLike} conclusion is unchanged.

There is a separate rank-one branch with polynomial \(X-1\).
That branch already supplies the stronger \(M\mid jx\), so it supplies the
weakened field by \(C\mid M\). It does not have to satisfy the degree-\(3^a\)
selected-root predicate.

The numbers \(2\) and \(3\) used in this lemma are \emph{order tests}.
They are not replacements for the two actual level primes \(p_0,p_1\).
Both actual level-prime blocks remain necessary: for any coefficient prime \(\ell\),
one of them is away from \(\ell\), and the proof takes its witness there.
The original \(M\) remains in the tame character and in the forbidden-root condition.

\section{The improved weighted projector estimate}
\label{sec:counting}

The exact source comparison is in
\src{New/Spanning/MainTerm/BadTypes}{New/Spanning/MainTerm/BadTypes.lean}.
For \(1<r<q\), every label contributes at least
\[
                         w_q=\frac{1-r/q}{q^r}
\]
to the coefficient at the identity of its idempotent. This coefficient is the
regular-representation dimension share. It is not the uniform probability
measure on all irreducible representations.

\subsection{Bound the complement, not just the witness mass}

Let \(B_q=|\labels_{q,r}|\). Among all \(q^r\) monic degree-\(r\) polynomials,
the uniform weight of the nonselected ones is
\begin{align}
    (q^r-B_q)w_q
       &\le \frac{q^r-B_q}{q^r}\notag\\
       &=1-\frac{B_q}{q^r}
        \le \exp\left(-\frac{B_q}{q^r}\right).
        \label{eq:complement}
\end{align}
Thus the useful exponent is
\begin{equation}
       u_q=\frac{B_q}{q^r}
          =\frac{1}{3r}
             \left(1-3q^{-2r/3}+2q^{-r}\right).
       \label{eq:uexact}
\end{equation}
There is no further loss of \(1-r/q\) in this exponent:
\eqref{eq:complement} bounds the weighted \emph{complement} from above,
exactly as the original \code{card_not_primitive_mul_le} does.

For \(A=q^{r/3}\ge2\),
\[
     A^3-3A+2\ge\tfrac12 A^3,
\]
because
\[
     A^3-6A+4=(A-2)(A^2+2A-2)\ge0.
\]
Consequently
\begin{equation}
                         u_q\ge\frac{1}{6r}.
                         \label{eq:usimple}
\end{equation}
The actual source has much larger \(A\); the simple bound deliberately
avoids optimizing an unimportant constant.

\subsection{The tuple argument}

Let there be \(f\) fields of cardinality \(q\), with the same rank \(r\).
The allowed-label counting lemma is retained:
if \(0\le\varepsilon\le1\) and
\((1+M+2r)/q\le\varepsilon\), the uniform weight of allowed labels in each
coordinate is at least \(1-\varepsilon\).

Define a new good tuple to have all entries allowed and at least one entry in
\(\labels_{q,r}\). The weighted count of all-allowed tuples is at least
\((1-\varepsilon)^f\). The weighted count of tuples with no selected entry is
at most \(e^{-fu_q}\). Subtracting this latter bound and using the lower
idempotent weight coordinate by coordinate gives
\[
        \operatorname{coeff}_1(e_{\mathrm{good}})
                 \ge(1-\varepsilon)^f-e^{-fu_q}.
\]
All coefficients in this comparison are real, or are read as real parts as in
the source. Bernoulli's inequality now yields
\begin{equation}
       \operatorname{coeff}_1(1-e_{\mathrm{good}})
          \le f\varepsilon+e^{-fu_q}
          \le f\varepsilon+e^{-f/(6r)}.
       \label{eq:bad}
\end{equation}

The selected labels are themselves allowed, since they are irreducible of degree
\(r>1\) and hence have no roots in the base field. The counting proof also works
without needing to exploit this inclusion: it is the same allowed-tuples-minus-
nonwitness-tuples cover already used by the source.

\subsection{Comparison with the inspected bound}

The original lemma counts primitive polynomials and applies
\[
                   n\le\varphi(n)(2\log n+1).
\]
It obtains exponent \(1/[2r(2r\log q+1)]\).
To force a one-block tail at most \(e^{-L}\), it suffices now that
\[
                         f\ge6rL,
\]
where previously the displayed sufficient condition was
\[
                         f\ge2r(2r\log q+1)L.
\]
The ratio of these sufficient thresholds is \((2r\log q+1)/3\).
For the source's \(L=9\log b\), the new local threshold is
\(f\ge54r\log b\).

The actual projector uses both level primes. If each block has \(f\) factors,
with errors \(\varepsilon_0,\varepsilon_1\), a combined bound is
\begin{equation}
       f(\varepsilon_0+\varepsilon_1)+2e^{-f/(6r)}.
       \label{eq:twoblock}
\end{equation}
The source budgets \(b^{-8}/2\) separately for each block before combining them,
so retaining that bookkeeping preserves its total \(b^{-8}\) allowance.
If one instead specifies a single combined exponential allowance \(e^{-L}\),
use \(f\ge6r(L+\log2)\).

The group-algebra estimate becomes an estimate on the actual automorphic vectors
through the source's unfolded moment calculation. It does not bound the lost
dimension in every arbitrary representation. That transfer remains part of
\src{New/Spanning/SecondMoment/SpanningTheorem}{SpanningTheorem.lean}.

\section{Concrete examples and failure modes}
\label{sec:examples}

The examples use exponent coordinates in the cyclic group \(\F_{q^r}^\times\).
If \(g\) is a generator and \(z=g^k\), then
\[
       \ord(z)=\frac{q^r-1}{\gcd(k,q^r-1)},\qquad
       \ord(\Norm(z))=\frac{q-1}{\gcd(k,q-1)}.
\]
These formulas make all the following computations exact.

\subsection{A strictly larger family}

Take \(q=37\), \(r=3\), and \(M=6\). Then
\[
       N=50652,\quad h=108,\quad K=469.
\]
The root \(g^7\) is selected: \(\gcd(7,6)=1\), and \(469\nmid7\).
It has order \(7236\), so it is not primitive, while its norm has order \(36\).
Its \(w=z^{36}\) has order \(201=3\cdot67\); removal of either prime-primary
component leaves the other obstruction to falling into the base field.
The old primitive-label count is \(4752\), whereas the new count is \(5616\).

\begin{center}
\begin{tabular}{@{}rrrrr@{}}
\toprule
\(q\) & \(r\) & \(M\) & Primitive labels & Selected labels\\
\midrule
37 & 3 & 6 & 4,752 & 5,616\\
73 & 3 & 6 & 43,200 & 43,200\\
109 & 3 & 6 & 123,120 & 143,856\\
1801 & 3 & 30 & 445,080,960 & 649,080,000\\
\bottomrule
\end{tabular}
\end{center}

The equality at \(q=73\) illustrates why strict enlargement should not be asserted
uniformly. The exact count and its simple lower bound are the useful uniform facts.
The last row is calculated by exact integer formulas, not exhaustive enumeration.

\subsection{Primitive norm does not suffice}

Continue with \(q=37\), \(r=3\), \(M=6\), but take \(z=g^{469}\).
Its norm has order \(36\), so it is primitive in \(\F_{37}^{\times}\).
The root itself has order \(108\) and degree \(3\). However,
\[
                        \ord(z^3)=36,
\]
so \(z^3\in\F_{37}\). The two order tests pass, but \(z^h=1\).
This is a counterexample under the full hypotheses \(6\mid M\), \(M^2\mid q-1\)
to dropping the subgroup exclusion.

\subsection{Full \texorpdfstring{\(M\)}{M}-order need not survive}

Take \(q=1801\), \(r=3\), \(M=30\). Then \(M^2=900\mid1800=q-1\).
The root \(g^{25}\) is selected and its norm has order \(72\).
That norm retains the full \(2\)-part \(8\) and the full \(3\)-part \(9\)
of \(1800\), but it has no factor \(5\).
This both demonstrates a larger family than the conservative condition retaining
every prime dividing \(M\) and explains why the internal order field
must actually be weakened.

\subsection{What the square divisibility is doing}

The assumption \(M^2\mid q-1\) makes multiplication by an \(M\)-torsion scalar
invisible to the two order tests. Merely \(M\mid q-1\) is insufficient in general.
For instance, in the exponent model at \(q=7,r=3,M=6\), adding the exponent
\(N/M=57\) of such a scalar changes \(k=1\) into \(58\), which is even and
fails the \(2\)-test. The source already has the needed square divisibility;
this proposal retains it.

\clearpage
\section{Implementation map}
\label{sec:integration}

The change should be implemented through a fresh predicate and a narrow API.
The new polynomials do not satisfy the old definition of primitive polynomial.
Keep \code{IsPrimitivePolynomial} as a useful existing notion; use a new
predicate, for example \code{IsTwoTestPolynomial}, for the enlarged family.
All names suggested here are proposed API names, not claims that declarations
with these names already exist in the repository.

\subsection{The six direct source files}

\begin{longtable}{@{}>{\raggedright\arraybackslash}p{.39\textwidth}>{\raggedright\arraybackslash}p{.56\textwidth}@{}}
\toprule
File under \code{src/} & Concrete change\\
\midrule
\endhead
\code{New/UpperBound/Objects/}\newline
\code{LabelsAndIndexTypes.lean}
& Introduce the selected-label predicate. Prove its label and irreducibility
projections, scalar stability, common norm power, and surviving \(2,3\)-order.
Replace the existential primitive witness in good tuples and good tuple
families. Preserve the allowed-label predicate.\\[5pt]
\code{New/Spanning/MainTerm/}\newline
\code{BadTypes.lean}
& Prove the count \eqref{eq:labelcount}, or its sufficient lower bound.
Replace \code{card_not_primitive_mul_le} in this branch by
\eqref{eq:complement}--\eqref{eq:usimple}. Generalize the tuple-count lemma
to accept a witness family with an exponential complement bound.\\[5pt]
\code{New/Spanning/Multiplicities/}\newline
\code{EqualMultiplicities.lean}
& Replace extraction of a primitive witness by extraction of a selected witness.
The local cuspidal-test argument uses only labelhood and irreducibility.\\[5pt]
\code{New/Congruences/}\newline
\code{GaloisRepresentations/}\newline
\code{SelmerHypotheses.lean}
& Replace the full-rank alternative in \code{exists_primitive_or_one}
by the selected-label predicate. Preserve the \(X-1\) rank-one alternative and
a witness above each actual level prime.\\[5pt]
\code{New/Congruences/}\newline
\code{GaloisRepresentations/}\newline
\code{AttachedRepresentations.lean}
& Fill the revised lift-setting field from the new good tuple families.
The root exclusions and the rank-one construction remain available.\\[5pt]
\code{New/Congruences/}\newline
\code{GaloisRepresentations/}\newline
\code{RowsAndColumnsFromLabels.lean}
& Build the first three common-root fields from Theorem~\ref{thm:main}.
Weaken its last field to \(6^{s+6}\mid jx\); use
Lemma~\ref{lem:nonresonance} at the unique scalar-order consumer.
Adapt the rank-one wrapper by divisibility.\\
\bottomrule
\end{longtable}

In \src{Lib/LocalGL/Green/LabelCharacters}{the Green-character file}, supplement the
generic lemma \code{injOn_Iio_pow_card_pow_of_isPrimitiveRoot}
with a weaker-input orbit lemma.
Its existing minimal-period conclusion is reusable; the replacement proof uses
\eqref{eq:w} instead of full primitive order.

The source anchors in \code{provenance/source_anchors.tsv} give exact starting lines,
Git blob identifiers, and SHA-256 hashes at the pinned revision. Text search is
evidence for locating these interfaces; it is not a substitute for a compiled
dependency check after the changes.

\subsection{A useful abstraction for the counting layer}

The tuple-count theorem is naturally independent of the witness's definition.
It can accept a predicate \(W\) on monic degree-\(r\) polynomials with:
\[
 W(c)\Longrightarrow\text{\(c\) is a label},\qquad
 \#\{c:\neg W(c)\}\,w_q\le e^{-u}.
\]
Then the same proof yields \(f\varepsilon+e^{-fu}\).
Instantiating this theorem with primitive polynomials should reproduce the old
result. Instantiating it with \(\labels_{q,r}\) gives the new result.
This separates the representation-theoretic bookkeeping from a particular
finite-field counting device and supplies a direct regression comparison.

For the arithmetic layer, retain a separate semantic theorem:
selected witnesses produce the common-root-power structure with the weakened
order field. Enlarging the projector is not justified by monotonicity alone.
Every added type must satisfy the Galois and congruence hypotheses, and that is
the purpose of the replacement theorem.

\subsection{A staged integration with clear stopping points}

\begin{enumerate}
\item Add finite-field lemmas and the selected-label predicate alongside existing
definitions. Check the inclusion of primitive labels and the small exact examples.
\item Generalize the weighted tuple lemma and prove the new bound while retaining
the original projector. At this stage the original final theorems should be untouched.
\item Generalize the common-root interface, adapt its sole scalar-order consumer,
and transfer the rank-one branch.
\item Change the good tuple families and lift-setting witness to the new predicate.
Rebuild the trace-comparison, Galois-representation, and congruence consumers.
\item Rebuild the final exported statements. Inspect their printed axiom dependencies
and the resolved import closure before describing the replacement as fully
kernel checked.
\end{enumerate}

This ordering makes failures local and gives each step a concrete comparison with
the existing development. It also permits keeping the old parameter choices at
first: the improved local bound fits their already established budgets.

\section{Verification, artifacts, and limitations}
\label{sec:verification}

\subsection{Exact executable checks}

The Python checks use only the standard library and perform exact integer
arithmetic in cyclic-group coordinates. For a group element \(g^k\), Frobenius
orbit length is computed independently as the multiplicative order of \(q\)
modulo \(\ord(g^k)\), rather than inferred from the proposed theorem.
For each tested selected root, the checker removes the entire primary component
at every prime dividing its order. Full degree after that maximal removal implies full
degree after any smaller power at the same prime.

The final fixed-two-test run reports:
\[
\begin{array}{lr}
\text{exponents inspected} & 1{,}754{,}908\\
\text{full-primary-removal checks} & 2{,}799{,}874\\
\text{scalar-scaling checks} & 3{,}610{,}932 .
\end{array}
\]
All checks passed.
The cases \((q,r,M)=(37,3,6),(73,3,6),(109,3,6)\) are exhaustive.
Larger deterministic samples cover
\((37,9,6),(433,3,12),(1801,3,30),(1297,3,36)\).
The last two explicitly test \(M\)-torsion scaling with extra prime factors or
higher \(2,3\)-powers. Sampled cases are labeled as sampled in the JSON;
their total exact counts come from formulas, not enumeration.

The scripts also record the examples and counterexamples from
Section~\ref{sec:examples}. An earlier, conservative checker retains the full
primary components at every prime dividing \(M\). It is included because the
final runner imports its elementary arithmetic helpers. Its stronger norm-order
conclusion must not be confused with the final two-test conclusion.

\subsection{Lean support}

The file \code{Lean/PowerOrbit.lean} is a standalone supplement using Lean
4.35.0 and the source project's pinned Mathlib revision
\code{56a005855ea17e0b0958f73ec801ae140b0a1679}.
It compiles with exit status \(0\), with no warnings.
The checked declarations and their axiom printouts are listed in
\code{Lean/lean_check.log}.

\begin{center}
\begin{tabularx}{\textwidth}{@{}lX@{}}
\toprule
Checked theorem & Formalized assertion\\
\midrule
\code{no_prime_power_kills} & A monoid element satisfying \eqref{eq:w},
with \(\gcd(3,K)=1\), cannot be killed by a prime-power.\\
\code{six_power_not_dvd} & For prime \(\ell\), \(0<jt<2^s\) implies
\(6^s\nmid j\ell^u t\).\\
\code{six_power_add_not_dvd} & The same nondivisibility for \(6^{s+c}\).\\
\code{power_orbit_injective} & The full \(3\)-power Frobenius-style orbit
of a root remains injective after an arbitrary prime-power operation,
under the explicit power hypotheses of the group lemma.\\
\bottomrule
\end{tabularx}
\end{center}

All four printed axiom lists contain only
\code{propext}, \code{Classical.choice}, and \code{Quot.sound};
none contains \code{sorryAx}.
The orbit proof uses Mathlib's minimal-period lemmas, as does the original
Green-character helper.

This file checks the delicate root-order and Frobenius-regularity ingredients,
plus scalar nonresonance. It is not a replacement for the six edited source
files. In particular, the complete finite-field polynomial count,
normalization API, weighted projector integration, and final RH/GRH import
closure have not been rebuilt with this change. The mathematical proofs in
this report cover those proposed bridges; their complete Lean integration
is a remaining implementation task.

\subsection{What this proposal establishes}

The defensible gain is an elementary replacement for primitive-root selection,
an exact label count, and an improved local projector estimate. It removes the
analytic totient lower bound from that counting branch and relaxes the local
witness requirement. The trace formula, theta-lift volume identities,
first-order congruence argument, Selmer bounds, and analytic zero criterion
remain necessary.

A reduction in total field size does not follow automatically. The sparse-place
construction separately requires a bound of the form
\[
                          2l_0(kt+1)<2^s,
\]
and several other admissibility inequalities remain. Reoptimizing all parameters
could be a subsequent task; no global rank, degree, line-count, or runtime
improvement is claimed here without doing that work.

The exact counts demonstrate that primitivity supplies unnecessary structure.
The sole-consumer order audit explains why discarding it is compatible with the
arithmetic mechanism. These two facts together make this a concrete simplification
candidate, with a small mathematical interface to validate before changing the
rest of the proof.

\appendix
\section{Reproducibility and source record}

The archive contains this \TeX{} source, its compiled PDF, a standalone Lean
supplement and log, two Python checkers and their JSON results, a count comparison
CSV, the proposal inventory, and the source-anchor manifest.
The top-level README gives commands to rebuild the PDF and rerun the checks.
A SHA-256 manifest covers the delivered artifacts.

The repository references throughout the report are immutable links to commit
\code{89d12dd40424f1217fd818d84877b3da52396a3b}.
The guide and the source are the primary evidence for descriptions of the supplied
proof. The new local assertions are proved in this report. The tests serve as
independent finite checks and regression artifacts, not as evidence that finite
enumeration proves a statement for all finite fields.

\begin{thebibliography}{9}
\bibitem{repo}
V. Reshetnikov, \emph{ProveIt: Generalized Riemann Hypothesis formalization},
\href{https://github.com/VladimirReshetnikov/ProveIt/tree/89d12dd40424f1217fd818d84877b3da52396a3b/Analysis/GeneralizedRiemannHypothesis/Lean/grh}{pinned Lean source tree},
commit \code{89d12dd40424f1217fd818d84877b3da52396a3b}.
\bibitem{guide}
\emph{Theta lifts and Dirichlet zeros}, \guide{}, same repository revision.
The local-label and counting sections explain the interfaces audited here.
\bibitem{proposals}
\href{https://github.com/VladimirReshetnikov/ProveIt/tree/89d12dd40424f1217fd818d84877b3da52396a3b/Analysis/GeneralizedRiemannHypothesis/docs/proposals}{Existing proposal collection}, same revision.
Individual paths and hashes appear in the accompanying inventory.
\bibitem{mathlib}
The Mathlib Community, \emph{Mathlib},
\href{https://github.com/leanprover-community/mathlib4/tree/56a005855ea17e0b0958f73ec801ae140b0a1679}{pinned revision}
\code{56a005855ea17e0b0958f73ec801ae140b0a1679},
with Lean 4.35.0. Used for the standalone formal checks.
\end{thebibliography}
\end{document}
